% from main.tex
%%%%%%%%%%  scan idx 542  |  folio 529  %%%%%%%%%%
alors de (1.5.2) et (1.5.1) :
\[
D_Y(x)-D_Y(y)=\dimdot\mathcal O_{S,y'}-\dimdot\mathcal O_{S,x'}+\degtr_{k(x')}k(x)-\degtr_{k(y')}k(y).
\]
Puisque $S$ est caténaire, on a
\[
\dimdot\mathcal O_{S,y'}-\dimdot\mathcal O_{S,x'}=\dimdot\mathcal O_{Z,y'}-\dimdot\mathcal O_{Z,x'}
\]
(EGA $0_{IV}$ 14.3.2.1), d'où
\begin{equation}\tag{1.5.4}
D_Y(x)-D_Y(y)=
[\dimdot\mathcal O_{Z,y'}-\degtr_{k(y')}k(y)]-[\dimdot\mathcal O_{Z,x'}-\degtr_{k(x')}k(x)].
\end{equation}
Puisque $Z$ est universellement caténaire, (EGA IV 5.6.5.1.) est applicable, et donne
\[
C_1=\dimdot\mathcal O_{Y,y}-e,\qquad C_2=\dimdot\mathcal O_{Y,x}-e,
\]
où $C_1$ et $C_2$ sont les deux termes à droite dans (1.5.4). Alors on conclut immédiatement de (1.5.4) que
\[
D_Y(x)-D_Y(y)=\dimdot\mathcal O_{Y,y}-\dimdot\mathcal O_{Y,x}>0,\qquad \text{C.Q.F.D.}
\]
On voit donc que (1.3.2) et (1.3.3) restent valables pour un schéma noethérien universellement caténaire $X$ de dimension $n$, si l'on définit la filtration de $K_\bullet(X)$ au moyen de (1.5.1) suivant la codimension des supports. On notera que cette filtration coincide avec celle de (1.1.1), lorsque de plus $X$ est biéqui-dimensionnel.

\subsection*{1.6}
On peut se demander s'il est toujours possible de définir sur un schéma noethérien $X$ de dimension finie une fonction $D$ possédant la propriété (i) de 1.4 et de plus la propriété suivante, plus forte que (ii) :
\begin{quote}
(ii bis) Si $x\in X$, et $y\ne x$ est une spécialisation immédiate de $x$, alors $D(x)=D(y)+1$.
\end{quote}
